% !TeX root = ../sections/03_solutions.tex
\subsection{対策-グラフ}
\begin{exercise}[段階関数のフーリエ係数とグラフ]
	周期 \(2\pi\) をもち，区間 \([-\pi,\pi)\) 上で
	\[
	f(x)=
	\begin{cases}
		0
		& \left(|x|\leq \dfrac{\pi}{3}\right),\\[6pt]
		\dfrac{\pi}{2}
		& \left(\dfrac{\pi}{3}<|x|\leq \dfrac{2\pi}{3}\right),\\[6pt]
		\pi
		& \left(\dfrac{2\pi}{3}<|x|\leq \pi\right)
	\end{cases}
	\]
	で与えられる周期関数を考える。
	
	\begin{enumerate}
		\item \(y=f(x)\) のグラフをかけ。
		\item \(f(x)\) のフーリエ係数を求めよ。
		\item gnuplot を用いて \(y=f(x)\)，\(y=S_7(x)\)，\(y=S_{31}(x)\) を
		一つの図にプロットせよ。
	\end{enumerate}
\end{exercise}

\begin{background}
	この関数は \(|x|\) によって定義されているため，偶関数である。
	すなわち，
	\[
	f(-x)=f(x)
	\]
	が成り立つ。
	
	したがって，フーリエ級数では正弦項が消える。
	つまり，
	\[
	b_n=0
	\]
	である。
	
	偶関数のフーリエ係数は，
	\[
	a_0=\frac{2}{\pi}\int_{0}^{\pi}f(x)\,dx
	\]
	および
	\[
	a_n=\frac{2}{\pi}\int_{0}^{\pi}f(x)\cos nx\,dx
	\]
	で求めればよい。
	
	この問題では，関数が区間ごとに定数なので，
	積分区間を
	\[
	0\leq x\leq \frac{\pi}{3},
	\qquad
	\frac{\pi}{3}<x\leq \frac{2\pi}{3},
	\qquad
	\frac{2\pi}{3}<x\leq \pi
	\]
	に分ける。
\end{background}
\begin{strategy}
	まずグラフは，原点を中心として左右対称な階段状のグラフである。
	
	\[
	|x|\leq \frac{\pi}{3}
	\]
	では高さ \(0\)，
	
	\[
	\frac{\pi}{3}<|x|\leq \frac{2\pi}{3}
	\]
	では高さ \(\dfrac{\pi}{2}\)，
	
	\[
	\frac{2\pi}{3}<|x|\leq \pi
	\]
	では高さ \(\pi\)
	である。
	
	次にフーリエ係数を求める。
	
	偶関数なので，
	\[
	b_n=0
	\]
	である。
	
	また，
	\[
	a_0
	=
	\frac{2}{\pi}
	\left(
	\int_{0}^{\pi/3}0\,dx
	+
	\int_{\pi/3}^{2\pi/3}\frac{\pi}{2}\,dx
	+
	\int_{2\pi/3}^{\pi}\pi\,dx
	\right).
	\]
	
	これを計算すると，
	\[
	a_0=\pi.
	\]
	
	次に \(a_n\) を求める。
	\[
	a_n
	=
	\frac{2}{\pi}
	\left(
	\int_{\pi/3}^{2\pi/3}\frac{\pi}{2}\cos nx\,dx
	+
	\int_{2\pi/3}^{\pi}\pi\cos nx\,dx
	\right).
	\]
	
	\[
	\int \cos nx\,dx=\frac{\sin nx}{n}
	\]
	を用いると，
	\[
	a_n
	=
	-\frac{1}{n}
	\left(
	\sin\frac{n\pi}{3}
	+
	\sin\frac{2n\pi}{3}
	\right).
	\]
	
	したがって，
	\[
	a_0=\pi,
	\qquad
	a_n
	=
	-\frac{1}{n}
	\left(
	\sin\frac{n\pi}{3}
	+
	\sin\frac{2n\pi}{3}
	\right),
	\qquad
	b_n=0.
	\]
	
	よって，第 \(N\) 部分和は
	\[
	S_N(x)
	=
	\frac{\pi}{2}
	+
	\sum_{n=1}^{N}
	a_n\cos nx
	\]
	である。
\end{strategy}
